\subsubsection*{REC01: Trusted resident models; utilization priority}
\begin{longtable}{>{\raggedright\arraybackslash}p{0.22\textwidth}>{\raggedright\arraybackslash}p{0.7\textwidth}}
\toprule
\textbf{Dimension} & \textbf{Assessment}\\\midrule\endhead
Candidate & Tuned MPS / software spatial controls\\\addlinespace
Evidence & P001 P002 P003 P008 E001-E004 E006 E031-E032\\\addlinespace
Expected benefit & Higher useful co-location under tested allocations\\\addlinespace
Sacrifice & Interference, quota tuning and memory pressure\\\addlinespace
Alternative / trigger & Partitions or dedication when trust/isolation dominates\\\addlinespace
Validation needed & Measure each model P99 and SLO goodput under bursts\\\addlinespace
\bottomrule\end{longtable}
\subsubsection*{REC02: Foreground tail-latency priority}
\begin{longtable}{>{\raggedright\arraybackslash}p{0.22\textwidth}>{\raggedright\arraybackslash}p{0.7\textwidth}}
\toprule
\textbf{Dimension} & \textbf{Assessment}\\\midrule\endhead
Candidate & Priority/preemptive scheduling\\\addlinespace
Evidence & P005 P006 P007 P015 P022 E018-E030 E044-E045 E052\\\addlinespace
Expected benefit & Lower foreground interference\\\addlinespace
Sacrifice & Background slowdown, profiling and runtime integration\\\addlinespace
Alternative / trigger & Dedicated GPU when supported control cannot meet deadline\\\addlinespace
Validation needed & Check foreground tails and background progress jointly\\\addlinespace
\bottomrule\end{longtable}
\subsubsection*{REC03: Low-rate serverless functions}
\begin{longtable}{>{\raggedright\arraybackslash}p{0.22\textwidth}>{\raggedright\arraybackslash}p{0.7\textwidth}}
\toprule
\textbf{Dimension} & \textbf{Assessment}\\\midrule\endhead
Candidate & Late binding / controlled runtime sharing\\\addlinespace
Evidence & P012 P013 E038-E041\\\addlinespace
Expected benefit & Higher function or workflow density\\\addlinespace
Sacrifice & Host CPU, memory transfer and shared-runtime trust\\\addlinespace
Alternative / trigger & Keep large or high-rate models resident if swaps dominate\\\addlinespace
Validation needed & Sweep model sizes, cold starts, offered rate and concurrency\\\addlinespace
\bottomrule\end{longtable}
\subsubsection*{REC04: Mixed inference with opportunistic training}
\begin{longtable}{>{\raggedright\arraybackslash}p{0.22\textwidth}>{\raggedright\arraybackslash}p{0.7\textwidth}}
\toprule
\textbf{Dimension} & \textbf{Assessment}\\\midrule\endhead
Candidate & Memory handover / controlled co-location\\\addlinespace
Evidence & P010 P019 E035 E048-E049\\\addlinespace
Expected benefit & Protect inference while using residual capacity\\\addlinespace
Sacrifice & Memory/reclamation costs and possible co-tenant regression; ArkGPU approximately 5\%QPS loss under low contention.\\\addlinespace
Alternative / trigger & Separate resources for intolerable interference\\\addlinespace
Validation needed & Measure inference tails separately from training throughput\\\addlinespace
\bottomrule\end{longtable}
\subsubsection*{REC05: Supported partition hardware; stable memory fit}
\begin{longtable}{>{\raggedright\arraybackslash}p{0.22\textwidth}>{\raggedright\arraybackslash}p{0.7\textwidth}}
\toprule
\textbf{Dimension} & \textbf{Assessment}\\\midrule\endhead
Candidate & MIG-aware hybrid policy\\\addlinespace
Evidence & P011 P014 P020 E036-E037 E042-E043 E050\\\addlinespace
Expected benefit & Resource allocation with configurable co-location\\\addlinespace
Sacrifice & Fragmentation and reconfiguration overhead\\\addlinespace
Alternative / trigger & Software sharing when partition granularity wastes capacity\\\addlinespace
Validation needed & Validate profile fit and online request tails; Dyna-P8.09\%makespan penalty at15\%prediction perturbation is local, not Section4.3.2simulation.\\\addlinespace
\bottomrule\end{longtable}
\subsubsection*{REC06: Several models with unequal service pressure}
\begin{longtable}{>{\raggedright\arraybackslash}p{0.22\textwidth}>{\raggedright\arraybackslash}p{0.7\textwidth}}
\toprule
\textbf{Dimension} & \textbf{Assessment}\\\midrule\endhead
Candidate & Fairness-aware admission / scheduling\\\addlinespace
Evidence & P018 E046-E047\\\addlinespace
Expected benefit & In \nolinkurl{Table3ratio3:1:2,total600requests/s:} goodput524.29 versusClockwork457.24; unfairness1.05 versus1.50.\\\addlinespace
Sacrifice & May give up aggregate throughput to improve fairness\\\addlinespace
Alternative / trigger & Priority scheduling when one foreground SLA dominates\\\addlinespace
Validation needed & Validate intended SLO/arrival mix; source uses100msSLO,50msperiod and max/min slowdown unfairness. No universal throughput advantage is inferred.\\\addlinespace
\bottomrule\end{longtable}
\subsubsection*{REC07: Predictable multi-model DNN co-location}
Abacus and LithOS broaden the candidate set where deterministic overlap or fine-grained resource control is compatible with the workload. Evaluate inference-only results separately from mixed training, preserve co-tenant sets, and measure profiling, host and kernel-support costs. Use a more conservative allocation when the protected service cannot tolerate the measured tail or unsupported kernels \cite{s021,s022}.

\subsubsection*{REC08: Edge deadlines with an explicit accuracy budget}
Consider Miriam for elastic-kernel coordination and Pantheon when chunk preemption and early-exit model variants are acceptable. The expected benefit is lower deadline interference under the evaluated mix; the sacrifice can include altered accuracy, extra resident variants, scheduling work and power. Prefer a model-preserving alternative when early exits violate task-quality requirements. Test nonperiodic arrivals, background progress and the actual platform rather than importing a single reported percentage \cite{s023,s025}.

\subsubsection*{REC09: Compute/bandwidth contention under a shared trust context}
SGDRC is a candidate for controlling compute and VRAM-bandwidth contention when its generated-kernel environment and shared CUDA context are acceptable. The expected benefit is latency-sensitive goodput with best-effort progress. It cannot be selected as a DNN runtime-fault boundary: that isolation is explicitly absent. Preserve cases where another baseline has greater raw throughput \cite{s024}.

\subsubsection*{REC10: Memory-constrained multi-LLM serving}
Oneiros is a candidate when a high-bandwidth CPU/GPU transfer path makes parameter remapping feasible. Its stated spatial-sharing case reports lower P99TBT/TTFT and higher tokens/s than vLLM. The sacrifice is transfer and remapping overhead; aggressive remapping can worsen average delay. Retain resident parameters or reduce remapping at low load or when transfers dominate, and establish required protection separately \cite{s026}.
\subsubsection*{REC11: Priority-aware CNN/BERT co-location with compilation access}
Consider ElasticRoom when generating kernel variants and maintaining offline tuning is feasible. Expect context-dependent real-time latency/Goodput improvement, while retaining model-specific no-gain cases, hardware drift and adverse generated kernels. Prefer a simpler or separately isolated mechanism when variant maintenance or a fault boundary is decisive \cite{s028}.
\subsubsection*{REC12: Static-DAG edge requests with mixed execution lengths}
Consider SPLIT when block-boundary preemption and its measured execution overhead fit the request budget. Short-request QoS can improve at the cost of long-request stability and spatial parallelism. Do not use the studied offline splitting method as evidence for input-dependent RNNs, arbitrary burst overload or universally reduced request-tail jitter \cite{s029}.
\paragraph{Known heterogeneous batch mixes.} Consider GACER when batch-completion efficiency is the objective and profiling/optimization delay is acceptable. Its reported 0.84--1.26 second optimization cost and absent serving-tail/fault evidence limit transfer to online SLO-sensitive or mutually untrusted tenants \cite{s027}.
